\documentclass[11pt]{article}

\usepackage{nmrdoc}
\usepackage{siunitx}
\usepackage{bm}

\sisetup{per-mode=symbol}
\DeclareSIUnit{\angstrom}{\text{\AA}}

\begin{document}
\NmrTitle{The McConnell Bond-Anisotropy Calculator}

\section{What the calculator computes}

The McConnell calculator describes the through-space magnetic effect of a
covalent bond whose electronic response is different along the bond axis than
across it. In a magnetic field, that anisotropic response creates a secondary
field at a nearby nucleus. The calculator represents the geometric part of that
effect as a rank-2 tensor: a \(3\times 3\) array whose \(ab\) entry says how a
field component applied along direction \(b\) contributes to the induced field
component along direction \(a\).

In the code, each accepted bond near each target atom contributes a tensor with
units \(\si{\angstrom^{-3}}\). The factor that would turn this geometric tensor
into a shielding contribution, such as a bond susceptibility anisotropy in the
usual McConnell formula, is not applied here. The output can therefore correlate
with a DFT shielding contribution after calibration, but before calibration it is
a geometric descriptor rather than a chemical shift or a shielding in ppm.

\section{Where shielding enters}

For a nucleus in an applied magnetic field \(\bm B_0\), the local field is
usually written
\begin{equation}
  \bm B_{\mathrm{nuc}} = (\bm I - \bm\sigma)\bm B_0 .
  \label{eq:shielding}
\end{equation}
Here \(\bm I\) is the \(3\times 3\) identity matrix and
\(\bm\sigma\) is the nuclear shielding tensor. The tensor is dimensionless in
physical NMR: it is the linear-response coefficient connecting the applied
field to the field screened by the electrons. The matrix form is needed because
a molecule is not a sphere. Rotating the molecule can change both the size and
the direction of the secondary field at the nucleus.

Equation~\eqref{eq:shielding} is a quantum-mechanical statement about electrons:
the applied field perturbs the electronic state, the perturbed electrons produce
their own magnetic field at the nucleus, and the shielding tensor records the
linear part of that response. The McConnell calculator keeps a classical
far-field approximation to one part of this response. It treats a bond as a
small anisotropic magnetic object located at its midpoint. The bond direction
sets the anisotropy axis, and the vector from the bond midpoint to the target
atom sets the observation direction. This is a local geometric model of the
angular response; it does not solve the electronic structure problem.

The NMR observable most often reported in solution is the isotropic chemical
shift. It is related to the isotropic shielding, the average of the three
diagonal tensor entries, after reference subtraction and sign conventions. That
scalar does not exhaust the calculation. The orientation-dependent part of
\(\bm\sigma\), carried by the symmetric traceless tensor components, is retained
because it records how the local geometry affects the angular pattern of the
shielding response.

\section{The tensor split used by the code}

A real \(3\times 3\) tensor \(\bm X\) is stored through a closed-form
split. The C++ type \texttt{SphericalTensor} owns this split. Its
\texttt{Decompose} routine is implemented in \texttt{Types.cpp}.
The scalar component is
\begin{equation}
  T_0 = \frac{1}{3}\operatorname{tr}\bm X .
  \label{eq:t0}
\end{equation}
This is the isotropic part: it multiplies the identity matrix and survives
orientational averaging.

The antisymmetric part is stored as a three-component pseudovector,
\begin{equation}
  \bm T_1 =
  \frac{1}{2}
  \begin{pmatrix}
    X_{yz}-X_{zy}\\
    X_{zx}-X_{xz}\\
    X_{xy}-X_{yx}
  \end{pmatrix}.
  \label{eq:t1}
\end{equation}
It is a compact way to store the matrix part that changes sign under transpose.
For standard isotropic solution NMR this part is usually not observed directly,
but the full calculator can produce it because the McConnell tensor is generally
not symmetric.

The remaining symmetric traceless matrix is
\begin{equation}
  \bm S =
  \frac{\bm X+\bm X^{\mathsf T}}{2}
  - T_0\bm I ,
  \qquad \operatorname{tr}\bm S = 0 .
  \label{eq:symmetric-traceless}
\end{equation}
It has five independent entries. The code packs them into the real spherical
\(T_2\) basis
\begin{equation}
  \bm T_2 =
  \begin{pmatrix}
    \sqrt{2}\,S_{xy}\\
    \sqrt{2}\,S_{yz}\\
    \sqrt{3/2}\,S_{zz}\\
    \sqrt{2}\,S_{xz}\\
    (S_{xx}-S_{yy})/\sqrt{2}
  \end{pmatrix}.
  \label{eq:t2}
\end{equation}
The factors in Eq.~\eqref{eq:t2} are the code's normalization. They preserve
the Frobenius norm: the squared length of this five-vector equals
\(\sum_{ab}S_{ab}^2\). The nine stored doubles are ordered as
\[
  [T_0,T_{1x},T_{1y},T_{1z},
    T_{2,-2},T_{2,-1},T_{2,0},T_{2,+1},T_{2,+2}] .
\]
For McConnell, the overall bond total is decomposed into the three parts. The
category tensor outputs are projected to the \(T_2\) subspace before they are
written.

\section{The method implemented}

For a target atom at position \(\bm r_i\) and a covalent bond \(j\) with endpoint
positions \(\bm q_A\) and \(\bm q_B\), the geometry module provides
\begin{equation}
  \bm m_j = \frac{\bm q_A+\bm q_B}{2}, \qquad
  \hat{\bm b}_j =
  \frac{\bm q_B-\bm q_A}{|\bm q_B-\bm q_A|}.
  \label{eq:bond-geometry}
\end{equation}
The midpoint \(\bm m_j\) has units of \(\si{\angstrom}\), while
\(\hat{\bm b}_j\) is a unitless bond-axis vector. Reversing the endpoint order
does not change the final tensor, because both appearances of
\(\hat{\bm b}_j\) change sign where needed.

The displacement from the bond midpoint to the target atom is
\begin{equation}
  \bm x_{ij} = \bm r_i-\bm m_j,\qquad
  r_{ij}=|\bm x_{ij}|,\qquad
  \hat{\bm d}_{ij}=\frac{\bm x_{ij}}{r_{ij}},\qquad
  c_{ij}=\hat{\bm d}_{ij}\cdot \hat{\bm b}_j .
  \label{eq:relative-geometry}
\end{equation}
The distance \(r_{ij}\) is in \(\si{\angstrom}\); \(c_{ij}\) is the cosine of
the angle between the bond axis and the midpoint-to-atom direction.

The full tensor contribution stored by the implementation is
\begin{equation}
  \mathcal M^{(ij)}_{ab} =
  \frac{
    9c_{ij}\,\hat d_{ij,a}\hat b_{j,b}
    -3\hat b_{j,a}\hat b_{j,b}
    -\left(3\hat d_{ij,a}\hat d_{ij,b}-\delta_{ab}\right)}
  {r_{ij}^{3}} .
  \label{eq:mcconnell-full}
\end{equation}
The Kronecker delta \(\delta_{ab}\) is one when \(a=b\) and zero otherwise.
The Cartesian indices \(a\) and \(b\) run over \(x,y,z\).
The numerator is dimensionless, so
\(\mathcal M^{(ij)}\) has units \(\si{\angstrom^{-3}}\). The first dyadic
product, \(\hat{\bm d}\otimes\hat{\bm b}\), is generally asymmetric and can
therefore produce \(T_1\). The second term is a projection onto the bond axis.
The third term is the negative of a symmetric traceless dipolar tensor.

For each accepted bond, the code also stores the dipolar tensor
\begin{equation}
  K^{(ij)}_{ab} =
  \frac{3\hat d_{ij,a}\hat d_{ij,b}-\delta_{ab}}{r_{ij}^{3}}
  \label{eq:dipolar}
\end{equation}
and the scalar
\begin{equation}
  f_{ij} =
  \frac{3c_{ij}^{2}-1}{r_{ij}^{3}} .
  \label{eq:scalar}
\end{equation}
Both \(K^{(ij)}\) and \(f_{ij}\) have units \(\si{\angstrom^{-3}}\).
The angular factor \(3c^2-1\) is the same two-lobed dipolar pattern carried by
the \(T_2\) tensor: it is positive near the bond axis, negative near directions
perpendicular to the bond, and zero at \(c=1/\sqrt{3}\), the magic angle of
about \(\ang{54.7}\). The scalar is the double contraction
\(\hat{\bm b}_j^{\mathsf T}\bm K^{(ij)}\hat{\bm b}_j\). For the full tensor in
Eq.~\eqref{eq:mcconnell-full}, \(\operatorname{tr}\mathcal M^{(ij)}=3f_{ij}\),
so \(f_{ij}\) is the \(T_0\) value of a single-bond full tensor, not the trace
of the dipolar tensor \(K\), whose trace is zero.

A scalar textbook McConnell term is often written with a prefactor such as
\((\Delta\chi/3)f_{ij}\), where \(\Delta\chi\) is the bond susceptibility
anisotropy. This implementation stops before that multiplication. It stores
\(\mathcal M^{(ij)}\), \(K^{(ij)}\), and \(f_{ij}\) as unscaled geometry, and it
does not choose separate susceptibility constants for C=O, C--N, aromatic, or
other bonds.

For each atom, the calculator queries bond midpoints within
\(\SI{10.0}{\angstrom}\). This is the McConnell bond-anisotropy cutoff in
\texttt{data/calculator\_params.toml}. A candidate contribution is rejected when
\(r_{ij}<\SI{0.1}{\angstrom}\), when the target atom is one of the bond
endpoints, or when \(r_{ij}\le 0.5\,\ell_j\), where \(\ell_j\) is the bond
length. The \(\SI{0.1}{\angstrom}\) guard prevents numerical divergence of the
\(r^{-3}\) expression. The \(0.5\,\ell_j\) rule excludes points inside the
near-field region of the finite bond source, where a point-source dipole
description is not the model being evaluated.

Accepted tensors are accumulated as
\begin{equation}
  \mathcal M_i = \sum_{j\in \mathcal A_i}\mathcal M^{(ij)},
  \label{eq:sum}
\end{equation}
where \(\mathcal A_i\) is the set of accepted nearby bonds for atom \(i\).
\(\mathcal M_i\) is decomposed by Eq.~\eqref{eq:t0}--\eqref{eq:t2} and written
as the nine-column full McConnell output. The source code field is named
\texttt{mc\_shielding\_contribution}; despite that name, the stored values are
the unscaled \(\si{\angstrom^{-3}}\) geometric tensor.

The calculator also keeps category summaries. The scalar \(f_{ij}\) is summed
separately for peptide C=O bonds, peptide C--N bonds, side-chain C=O bonds, and
aromatic bonds. Full tensors are summed into three broader groups: backbone
bonds, side-chain bonds, and aromatic bonds. Before writing those group tensors
as category features, the code symmetrizes each group tensor and subtracts its
trace divided by three, so those outputs contain the five \(T_2\) components
and no \(T_0\) or \(T_1\) part. The nearest peptide C=O and nearest peptide
C--N bonds are selected by midpoint distance; their stored \(T_2\) descriptors
come from the dipolar
tensor \(K\) in Eq.~\eqref{eq:dipolar}. Disulfide and unclassified bonds still
enter the overall total after passing the filters, but they do not enter the
written category \(T_2\) groups or scalar sums.

The main approximations are therefore visible in the code path. A bond is
collapsed to its midpoint and unit axis. Accepted bonds have unit geometric
weight in the base calculator; no bond-order, element, or susceptibility
coefficient is multiplied into Eq.~\eqref{eq:mcconnell-full}. The \(r^{-3}\)
form is used outside a fixed near-field exclusion, and the category labels are
used for output grouping rather than for changing the formula.

The same tensor expression can also be sampled at an arbitrary point in space.
That path iterates over the bond list, accepts bonds within
\(\SI{10.0}{\angstrom}\), applies the \(\SI{0.1}{\angstrom}\) singularity
guard, and omits points with
\(r_{ij}<0.5\,\ell_j\). The per-atom path requires \(r_{ij}>0.5\,\ell_j\), so
it rejects equality at that near-field boundary as well. The free-point path
does not run the endpoint self-source
test, because a sampling point has no atom index. When the nearest peptide C=O
direction is stored on the atom, the code re-normalizes it unless its norm is
below \(10^{-10}\), in which case the direction is written as zero.

\section{Inputs and output}

This calculator consumes molecular geometry: atom positions in
\(\si{\angstrom}\), typed covalent bonds, bond categories, bond midpoints, bond
lengths, and bond directions. Its declared runtime dependencies are the geometry
result and the spatial index over bond midpoints. It does not consume MOPAC
charges, AIMNet2 charges or embeddings, ff14SB charges, APBS electrostatic
fields, or DSSP secondary-structure data. Those inputs feed other charge,
electrostatic, and secondary-structure calculations in the pipeline. MOPAC bond
orders are used by the separate MOPAC-weighted McConnell variant, not by this
calculator.

The feature files reflect that geometry-derived role. \texttt{mc\_shielding.npy}
stores an \(N\times 9\) table: the nine \(T_0+T_1+T_2\) components of
\(\mathcal M_i\) for each atom. \texttt{mc\_category\_T2.npy} stores an
\(N\times 25\) table with five five-component \(T_2\) blocks: backbone total,
side-chain total, aromatic total, nearest peptide C=O, and nearest peptide
C--N. \texttt{mc\_scalars.npy} stores an \(N\times 6\) table with the four scalar sums
\((f_{\mathrm{CO}},f_{\mathrm{CN}},f_{\mathrm{sidechain\,CO}},f_{\mathrm{arom}})\)
and the nearest peptide C=O and peptide C--N midpoint distances. If no accepted
bond of one of those nearest-bond types is found, the stored distance remains
\(\SI{99.0}{\angstrom}\), the code's missing-value marker. Tensor and scalar
values remain in \(\si{\angstrom^{-3}}\). Calibration against a DFT shielding
reference is the later step that can attach physical scale and compare these
descriptors with shielding changes.

\end{document}
